\mode*
\section{Expected results}
\label{results}

% XXX(#5) This section outlines the expected results; to be replaced with the
% actual results once the prgi26 data has been collected and analysed.

The study is planned for the autumn 2026 instance of the course; here we
outline what the hypotheses predict that we will see in the data.

\begin{frame}
  \begin{description}
    \item[\Cref{h:patterns}] The debugging episodes that end in a fix contain contrast,
      generalisation and fusion; episodes that stall or end in
      trial-and-error \enquote{fixes} lack one or more patterns---we
      particularly expect to see generalisation without contrast (unguided
      experimenting) and contrast without generalisation (a single lucky or
      unlucky comparison, never validated).
    \item[\Cref{h:deep}] The number and completeness of patterns per episode correlates
      positively with the R-SPQ-2F deep-approach score, and negatively (or
      not at all) with the surface-approach score.
    \item[\Cref{h:success}] Fixing the bug, and doing so in fewer edit--run cycles,
      correlates positively with both the deep-approach score and pattern
      completeness---and if the overarching hypothesis is right, pattern
      completeness should \emph{mediate} the association between deep
      approach and debugging success.
  \end{description}
\end{frame}

For \cref{rq:experts}, the pilot (\cref{pilot})
leads us to expect the teachers' episodes
to exhibit all three patterns, with the explanation arriving only once the
patterns are complete; the interesting result is where the students'
episodes deviate from this expert baseline---which patterns are missing,
and in which order the remaining ones come.

For \cref{rq:generative}, \cref{h:generative} predicts that the students
who use the course material generate more, and more complete, patterns of
variation than the students who study on their own---with the difference
persisting after controlling for deep-approach score and prior programming
experience.
If instead the self-studiers generate just as many patterns but score higher
on the deep approach, that would suggest the ability comes with the approach
rather than with the instruction---also an informative outcome.

For \cref{rq:school}, \cref{h:school} predicts that students whose
gymnasiet had an independent principal score somewhat higher on the
surface-approach scale, and turn up more often among the students with the
weakest debugging test scores, smallest gains and least complete patterns.
We expect the association to be small if it is there at all, and we expect
to report it as a distribution rather than as an estimate: the huvudman
type is a weak covariate, the cohort is one course, and the schools were
not assigned (\cref{school-lookup}).
A flat result is the outcome we would bet on, and it would leave the
supply-side explanation of \cref{schooling} standing unopposed.

Beyond the hypotheses, the coded episodes will give us a catalogue of the
students' actual debugging behaviour---which patterns novices generate
spontaneously, in which order, and where they get stuck.
This catalogue is valuable in itself for designing debugging instruction: the
missing patterns are exactly what the instruction needs to provide.
